\chapter{Le décor : la pile complète}

Avant d'écrire une seule ligne de code, il faut savoir \emph{qui fait quoi}. Ce
chapitre pose la carte du terrain. Il se lit de bas en haut : de la fenêtre que
le système d'exploitation nous donne, jusqu'au bouton sur lequel l'utilisateur
clique. À la fin, vous saurez nommer chaque étage, dire de quoi il dépend, et —
c'est le point le plus important du chapitre — vous comprendrez pourquoi
l'étage du haut ne connaît pas l'étage du bas.

\section{Vue d'ensemble}

\begin{nkcode}{Vue d'ensemble — synthèse (voir Documentation, section « Raccordement au reste du moteur »)}
   Application  (NKGuiDemo, NKCode, NK3DModeler…)
        |  cree une NkWindow, pompe les NKEvent -> remplit ctx.input
        |  declare ses widgets entre BeginFrame / EndFrame
        v
   NKGui  (NkGuiContext -> NkGuiDrawList : vtx + idx + cmds)
        |  ne connait AUCUN GPU
        v
   Backend (au choix)
     |- NkGuiCanvasBackend   (header-only, dans NKCanvas/UI/)  -> NkIRenderer2D
     `- NkGuiRHIBackend      (Integrations/NKGui/)             -> NKRHI
        v
   NKCanvas (NkBatchRenderer2D -> SubmitBatches)  |  NKRHI
        v
   GPU
\end{nkcode}

Cinq étages, donc, mais qui ne s'empilent pas comme on l'imaginerait. Prenons-les
un par un.

\section{Étage 1 — la fenêtre : \texttt{NKWindow}}

\texttt{NKWindow} est l'abstraction de fenêtre native. C'est le seul module de la
pile qui parle au système d'exploitation : \texttt{HWND} sous Windows,
\texttt{Window} X11 ou surface Wayland sous Linux, \texttt{NSWindow} sous macOS.

Son usage se réduit à trois gestes :

\begin{nkcode}{Applications/NKGuiDemo/src/NKGuiDemo/main.cpp:243-251}
    NkWindow window;
    NkWindowConfig cfg;
    cfg.title = "NKGui - Demo Phase 2 (Coeur : drawlist + interaction)";
    cfg.width = 1100;
    cfg.height = 800;
    cfg.centered = true;
    cfg.resizable = true;
    if (!window.Create(cfg))
        return -1;
\end{nkcode}

\noindent puis, dans la boucle, \texttt{window.IsOpen()} pour savoir si elle vit
encore, et \texttt{window.SetCursor(\dots)} pour changer le pointeur. Le reste —
taille courante, position, état maximisé — s'interroge par des accesseurs
(\texttt{GetSize()} renvoie un \texttt{math::NkVec2u}).

Une option mérite d'être signalée dès maintenant parce qu'elle explique
l'apparence des applications du dépôt :

\begin{nkcode}{Engine/NKEditorKit/src/NKEditorKit/NkEditorShell.cpp:234}
    wc.frame = false;          // SANS bordure OS -> barre de titre custom (VSCode)
\end{nkcode}

Avec \texttt{frame = false}, le système ne dessine plus ni barre de titre ni
bordure : l'application les dessine elle-même, comme le fait Visual Studio Code.
C'est ce que fait \texttt{NKCode}. Notre application du chapitre~4 gardera la
décoration native, qui est plus simple.

\begin{nkpiege}[X11 définit \texttt{None}]
Sous Linux, \texttt{NKWindow} tire \texttt{<X11/Xlib.h>}, qui contient
\texttt{\#define None 0L}. Or plusieurs énumérations de NKGui ont un membre
nommé \texttt{None}. Le préprocesseur le transforme alors en \texttt{0L = 0} et
le compilateur signale « expected identifier » sur des lignes parfaitement
correctes, très loin de la vraie cause. La parade est en tête de l'en-tête
parapluie de NKGui :

\begin{nkcode}{Kernel/Runtime/NKGui/src/NKGui/NKGui.h:18-27 (abrégé)}
// Retire les macros X11 (None, Bool, Status...) AVANT toute declaration de
// NKGui. Sur Linux, NKWindow tire <X11/Xlib.h>, dont le `#define None 0L`
// transforme ensuite `None = 0` — membre de plusieurs enumerations de NkGuiTypes
// — en `0L = 0`. Le compilateur signale alors « expected identifier » sur des
// lignes parfaitement correctes, tres loin de la vraie cause.
#include "NKPlatform/NkX11Clean.h"
\end{nkcode}

La leçon générale : \textbf{incluez toujours \nkfile{NKGui/NKGui.h} et jamais un
en-tête interne de NKGui}. L'ordre d'inclusion de l'en-tête parapluie n'est pas
arbitraire, c'est l'ordre des dépendances.
\end{nkpiege}

\section{Étage 2 — les événements : \texttt{NKEvent}}

\texttt{NKEvent} transforme les messages du système en \textbf{événements
typés} : \texttt{NkWindowCloseEvent}, \texttt{NkMouseMoveEvent},
\texttt{NkMouseButtonPressEvent}, \texttt{NkKeyPressEvent},
\texttt{NkTextInputEvent}\dots

Le modèle est celui des \emph{callbacks} enregistrés une fois pour toutes, puis
d'une file qu'on vide à chaque image :

\begin{nkcode}{Applications/NKGuiDemo/src/NKGuiDemo/main.cpp:326-331 (extrait)}
    auto &events = NkEvents();
    bool running = true;
    events.AddEventCallback<NkWindowCloseEvent>([&](NkWindowCloseEvent *) { running = false; });
    events.AddEventCallback<NkMouseMoveEvent>([&](NkMouseMoveEvent *e) {
        ctx.input.mousePos = {static_cast<float32>(e->GetX()), static_cast<float32>(e->GetY())};
    });
\end{nkcode}

et, dans la boucle :

\begin{nkcode}{Applications/NKGuiDemo/src/NKGuiDemo/main.cpp:438-440}
        while (NkEvent *ev = NkEvents().PollEvent()) {
            (void)ev;
        }
\end{nkcode}

Ce \texttt{while} déroutant mérite un mot : \textbf{il ne fait rien du résultat}.
Le travail a déjà été accompli par les \emph{callbacks} enregistrés plus haut ;
\texttt{PollEvent} sert uniquement à \emph{pomper} la file jusqu'à ce qu'elle
soit vide, ce qui déclenche au passage l'appel de tous les \emph{callbacks}. On
retrouve exactement le même idiome dans la boucle principale de
\texttt{NKEditorKit} (\nkfile{Engine/NKEditorKit/src/NKEditorKit/NkEditorShell.cpp:314}).

Une distinction fondamentale, qu'on retrouvera au chapitre~4 :

\begin{itemize}
  \item \texttt{NkTextInputEvent} porte le \textbf{texte} tapé, déjà décodé par
        le système en \emph{codepoint} Unicode — claviers AZERTY et méthodes de
        saisie asiatiques comprises. C'est lui qui alimente les champs de
        saisie ;
  \item \texttt{NkKeyPressEvent} / \texttt{NkKeyReleaseEvent} portent les
        \textbf{touches} : flèches, Suppr, Entrée, F2, Ctrl+quelque chose. Ce
        sont elles qui alimentent la navigation.
\end{itemize}

Confondre les deux est un classique : on branche les flèches sur le texte, et
plus rien ne fonctionne. Le commentaire du pont de NKEditorKit le dit :

\begin{nkcode}{Engine/NKEditorKit/src/NKEditorKit/NkEditorShell.cpp:367-370}
// Mappe une touche OS d'edition vers l'etat ENFONCE NKGui (press/release ->
// repetition au maintien geree par NKGui). Sans ce pont, aucune navigation
// clavier ni edition de texte dans les panneaux.
\end{nkcode}

\section{Étage 3 — le GPU. Attention, il y a deux abstractions}

C'est ici que le débutant se perd, et c'est normal : le moteur contient
\textbf{deux} abstractions GPU distinctes, qui ne se superposent pas.

\subsection{\texttt{NKRHI} — l'abstraction GPU bas niveau}

\nkfile{Kernel/Runtime/NKRHI} est l'abstraction « moderne » du GPU :
périphérique (\texttt{NkIDevice}), tampons de commandes, passes de rendu,
\emph{pipelines}, descripteurs. Son arborescence de backends parle d'elle-même :
\texttt{DirectX11}, \texttt{DirectX12}, \texttt{Metal}, \texttt{Opengl},
\texttt{Software}, \texttt{Vulkan}. C'est sur elle que repose le rendu 3D
(\texttt{NKRenderer}) et, par ricochet, les applications qui font de la 3D.

\subsection{\texttt{NKCanvas} — sa propre abstraction, pour la 2D}

\texttt{NKCanvas} \textbf{ne passe pas par NKRHI}. Il possède ses propres
contextes graphiques et ses propres renderers 2D. La preuve tient en une ligne
de son fichier de build :

\begin{nkcode}{Kernel/Runtime/NKCanvas/NKCanvas.jenga:65}
    _canvasDeps = ["NKWindow", "NKFont", "NKImage", "NKStream", "NKTime", "NKGlad", "NKThreading"]
\end{nkcode}

Pas de \texttt{NKRHI} dans la liste. L'arborescence
\nkfile{Kernel/Runtime/NKCanvas/src/NKCanvas/Backend/} contient
\texttt{DirectX/} (DX11 et DX12), \texttt{Metal/}, \texttt{OpenGL/},
\texttt{Software/} et \texttt{Vulkan/} : ce sont les implémentations \emph{de
NKCanvas}, pas celles de NKRHI.

Il y a donc, dans le dépôt, deux familles de backends portant les mêmes noms
d'API graphique et ne se connaissant pas. Cette duplication est assumée et même
documentée à l'endroit où elle pose problème :

\begin{nkcode}{Engine/NKEditorKit/src/NKEditorKit/NkIEditorRenderer.h:28-31}
// Choix d'API graphique NEUTRE (decouple de NKCANVAS ET NKRHI : leurs enums
// NkGraphicsApi se dupliquent dans le namespace nkentseu et ne peuvent
// cohabiter dans un meme TU). Chaque impl mappe vers son propre enum.
enum class NkEditorGfxApi : uint8 { Auto = 0, OpenGL, Vulkan, DX11, DX12, Software };
\end{nkcode}

\noindent Lisez bien : les deux énumérations \texttt{NkGraphicsApi} — celle de
NKCanvas et celle de NKRHI — \emph{ne peuvent pas cohabiter dans une même unité
de traduction}. D'où le besoin d'un troisième énuméré, neutre, quand une couche
doit pouvoir parler aux deux.

\begin{nkretenir}[Une fenêtre = une pile]
C'est la doctrine du dépôt, énoncée mot pour mot :

\begin{nkcode}{Integrations/NKGui/NkEditorRHIRenderer.h:8-13 (extrait)}
* toute application moteur (NkAnimaEditor, Nogee, ...) qui veut la coquille
* NkEditorShell rendue sur NKRHI/NKRenderer (et PAS NKCanvas — regle
* « une fenetre = une pile ») injecte cette classe via
* NkEditorShellConfig::renderer.
\end{nkcode}

Une fenêtre est rendue \textbf{soit} par NKCanvas, \textbf{soit} par NKRHI —
jamais les deux. Le choix se fait une fois, à la création. Pour ce cours, ce sera
toujours NKCanvas.
\end{nkretenir}

\subsection{Les cinq backends 2D de NKCanvas, et leur état réel}

Les cinq dossiers de backends ne fournissent pas tous un renderer 2D, et ceux
qui le fournissent ne sont pas tous validés à l'exécution. Voici l'état, tel
qu'il ressort des sources et des feuilles de route :

\begin{center}
\begin{tabular}{@{}lll@{}}
\toprule
\textbf{Backend} & \textbf{Renderer 2D} & \textbf{État à l'exécution} \\
\midrule
OpenGL   & \texttt{NkOpenGLRenderer2D} (926 l.)  & le seul \textbf{validé} \\
Vulkan   & \texttt{NkVulkanRenderer2D} (1524 l.) & le code le plus complet ; plantages signalés \\
DX12     & \texttt{NkDX12Renderer2D} (1320 l.)   & plantages signalés \\
DX11     & \texttt{NkDX11Renderer2D} (764 l.)    & écran vide/noir signalé \\
Software & \texttt{NkSoftwareRenderer2D} (674 l.) & écran vide/noir signalé \\
Metal    & \textbf{aucun}                        & non implémenté \\
\bottomrule
\end{tabular}
\end{center}

La ligne « Metal » n'est pas une supposition, elle est écrite dans la fabrique :

\begin{nkcode}{Kernel/Runtime/NKCanvas/src/NKCanvas/Renderer/Core/NkRenderer2DFactory.cpp:86-89}
case NkGraphicsApi::NK_GFX_API_METAL:
    NK_R2D_FACTORY_ERR("Metal 2D renderer not yet implemented");
    return nullptr;
\end{nkcode}

Quant aux autres verdicts, ils viennent de
\nkfile{Kernel/Runtime/NKCanvas/ROADMAP.md} (lignes 185-200 et 601-607), à la
date du 30 mai 2026 : seul OpenGL était alors validé à l'exécution sur la démo
Pong. \textbf{Ces états sont datés et méritent d'être revérifiés} avant de tirer
une conclusion définitive : le dépôt bouge. Mais ils expliquent pourquoi, quand
une application « ne montre rien », la première chose à essayer est de changer de
backend.

Le choix, lui, est explicite. Soit on nomme l'API :

\begin{nkcode}{Applications/NKGuiDemo/src/NKGuiDemo/main.cpp:253-262}
    NkContextDesc desc;
    desc.api = NkGraphicsApi::NK_GFX_API_AUTO;
    if (desc.api == NkGraphicsApi::NK_GFX_API_AUTO) {
#if defined(NKENTSEU_PLATFORM_WINDOWS)
        desc.api = NkGraphicsApi::NK_GFX_API_DX11;
#else
        desc.api = NkGraphicsApi::NK_GFX_API_OPENGL;
#endif
    }
\end{nkcode}

\noindent soit on laisse la fabrique essayer une liste, par
\texttt{NkContextFactory::CreateWithFallback(window, preferenceOrder, count)}.
L'ordre conseillé documenté en \nkfile{.../Factory/NkContextFactory.h:46} est
\texttt{\{DX12, DX11, Metal, Vulkan, OpenGL, Software\}}.

\section{Étage 4 — \texttt{NkCanvas}, le rendu 2D}

\texttt{NkCanvas} est décrit dans \nkfile{Guides/05-NKCanvas.md} comme
« la couche de rendu 2D conviviale de Nkentseu, l'équivalent de la partie
\emph{Graphics} de SFML ». Il fournit :

\begin{itemize}
  \item un \textbf{contexte graphique} (\texttt{NkIGraphicsContext}) créé par une
        fabrique selon l'API choisie ;
  \item une \textbf{cible de rendu} (\texttt{NkRenderTarget}), dont la
        spécialisation \texttt{NkRenderWindow} est celle qu'on utilise
        99\,\% du temps : elle possède le cycle d'image ;
  \item un \textbf{renderer 2D} (\texttt{NkIRenderer2D}, façade
        \texttt{NkRenderer2D}) qui offre les primitives : lignes, rectangles,
        cercles, triangles, et le point d'entrée bas niveau
        \texttt{DrawVertices} ;
  \item des \textbf{ressources} : \texttt{NkTexture}, \texttt{NkFont},
        \texttt{NkSprite}, \texttt{NkText}, \texttt{NkShader} ;
  \item des \textbf{formes persistantes} façon SFML : \texttt{NkRectangleShape},
        \texttt{NkCircleShape}, \texttt{NkLineShape}, \texttt{NkConvexShape}.
\end{itemize}

Le chapitre~2 lui est entièrement consacré. Retenez pour l'instant une seule
chose : NkCanvas dessine des \textbf{triangles}. Tout le reste — cercles,
lignes épaisses, texte — est fabriqué à partir de triangles par le module
lui-même.

\section{Étage 5 — \texttt{NKGui}, les widgets}

\texttt{NKGui} est le framework d'interface : boutons, cases à cocher, sliders,
champs de saisie, arbres, tables, menus, fenêtres flottantes, docking. Environ
7\,780 lignes réparties sur 13 fichiers source, dont
\nkfile{Kernel/Runtime/NKGui/src/NKGui/Widgets/NkGuiWidgets.cpp} qui pèse à lui
seul 4\,973 lignes.

Son arborescence est petite et lisible — c'est par là qu'il faut commencer :

\begin{nkcode}{Kernel/Runtime/NKGui/ — arborescence}
Kernel/Runtime/NKGui/
  NKGui.jenga · ARCHITECTURE.md · README.md · ROADMAP.md
  src/NKGui/
    NKGui.h                     <- en-tete PARAPLUIE (la seule chose a inclure)
    NkGuiExport.h -> NkGuiApi.h <- macros d'export (defaut : statique)
    Core/
      NkGuiTypes.h      (333 l.) types de base : NkGuiId, enums, themes de flags
      NkGuiInput.h      (145 l.) etat d'entree par frame (header-only, tout inline)
      NkGuiDrawList.h/.cpp (101/342 l.) la LISTE DE COMMANDES de dessin
      NkGuiFont.h/.cpp  ( 81/158 l.) police + atlas (wrapper NKFont)
      NkGuiContext.h/.cpp (551/573 l.) LE contexte : etat complet de l'UI
    Widgets/
      NkGuiWidgets.h/.cpp (405/4973 l.) TOUS les widgets
\end{nkcode}

\subsection{Le point capital : NKGui ne connaît aucun GPU}

Voici la ligne la plus importante de ce chapitre :

\begin{nkcode}{Kernel/Runtime/NKGui/NKGui.jenga:22-27}
nkentseudependson(
    ["NKPlatform", "NKCore", "NKMemory", "NKMath", "NKThreading",
     "NKLogger", "NKContainers", "NKEvent", "NKFont", "NKImage"],
    selfexport="NKGui",
    extra_includes=["src"],
)
files(["src/**.cpp"])
\end{nkcode}

Relisez la liste. Il n'y a \textbf{ni \texttt{NKCanvas}, ni \texttt{NKRHI}, ni
même \texttt{NKWindow}}. NKGui ne connaît que les mathématiques, les conteneurs,
la mémoire, les polices, les images et les événements. Il ne sait pas ce qu'est
une texture GPU, ni un \emph{shader}, ni une fenêtre.

\begin{nkretenir}
\textbf{NKGui ne dépend ni de NkCanvas ni de NKRHI.} Il ne dessine rien : il
\emph{produit une liste de commandes} — des sommets, des indices, des commandes
de dessin — et laisse quelqu'un d'autre les envoyer au GPU. Ce quelqu'un
s'appelle un \textbf{backend}, et le mariage se fait par un \textbf{pont}.
\end{nkretenir}

\subsection{Ce que NKGui produit : \texttt{NkGuiDrawList}}

La sortie de NKGui est purement géométrique :

\begin{nkcode}{Kernel/Runtime/NKGui/src/NKGui/Core/NkGuiDrawList.h:23-46 (abrégé)}
struct NkGuiVertex {
        NkVec2 pos;
        NkVec2 uv;   ///< (0,0) = couleur unie (pas de texture)
        uint32 col;
};

enum class NkGuiDrawCmdType : uint8 {
    Triangles,         ///< triangles unis
    TexturedTriangles  ///< triangles textures (texte/atlas/image)
};

struct NkGuiDrawCmd {
        NkGuiDrawCmdType type = NkGuiDrawCmdType::Triangles;
        uint32 idxOffset = 0;
        uint32 idxCount  = 0;
        uint32 texId     = 0;
        NkRect clipRect  = {0.f, 0.f, 1.0e9f, 1.0e9f};   // 1e9 = « pas de clip »
};
\end{nkcode}

\begin{nkcode}{Kernel/Runtime/NKGui/src/NKGui/Core/NkGuiDrawList.h:48-56 (abrégé)}
struct NkGuiDrawList {
        NkVector<NkGuiVertex> vtx;
        NkVector<uint32>      idx;
        NkVector<NkGuiDrawCmd> cmds;
        NkRect  clipStack[32] = {};
        int32   clipDepth = 0;
        float32 thickScale = 1.f; ///< echelle DPI des epaisseurs (preservee par Reset)
};
\end{nkcode}

Trois tableaux, donc : les sommets, les indices, et les commandes. Une
\emph{commande} décrit une tranche d'indices à dessiner avec une texture donnée
et un rectangle de découpe donné. Notez le \texttt{texId} : c'est un simple
entier, pas un pointeur vers une texture GPU. NKGui manipule des
\emph{identifiants}, à charge du backend de savoir à quoi ils correspondent.

Notez aussi la sentinelle \texttt{1e9} pour « pas de découpe ». Le backend la
reconnaît par un test tolérant (\texttt{w < 1e8}), ce qui évite toute
comparaison exacte de flottants.

\section{Le pont : marier NKGui et NkCanvas}

Puisque NKGui ne connaît pas NkCanvas, et que NkCanvas ne connaît pas NKGui, il
faut un tiers qui connaisse les deux. Le dépôt en fournit deux, selon le public.

\subsection{\texttt{NkGuiCanvasBackend} — le pont vers NkCanvas}

Il vit dans \nkfile{Kernel/Runtime/NKCanvas/src/NKCanvas/UI/NkGuiCanvasBackend.h}
et il est \textbf{header-only}. Ce détail est délibéré et sa justification est en
tête de fichier :

\begin{nkcode}{Kernel/Runtime/NKCanvas/src/NKCanvas/UI/NkGuiCanvasBackend.h:2-9}
// NkGuiCanvasBackend.h — rend un nkgui::NkGuiDrawList via NKCanvas (NkIRenderer2D).
// Backend RÉUTILISABLE (lib) : le cœur NKGui reste render-agnostique ; ce pont
// traduit ses draw-lists en appels NkIRenderer2D. Gère l'atlas de police
// (gray8 -> RGBA blanc+alpha) et les images RGBA pour les commandes texturées.
//
// HEADER-ONLY : NKCanvas ne le compile pas (pas de .cpp) ; seuls les consommateurs
// (qui dépendent déjà de NKGui ET NKCanvas) l'incluent. Modelé sur NkUICanvasBackend.
\end{nkcode}

Comprenez bien l'astuce. Ce fichier inclut \texttt{"NKGui/NKGui.h"}, donc il
dépend de NKGui. Mais il vit dans NKCanvas, dont le \texttt{files([...])} ne
prend que les \texttt{.cpp} : \textbf{personne ne le compile dans NKCanvas}. Il
n'est compilé que dans l'application qui l'inclut — application qui, elle,
dépend bien des deux modules. Résultat : NKCanvas n'acquiert aucune dépendance
vers NKGui, et le pont existe quand même.

Son API tient en quatre méthodes :

\begin{nkcode}{Kernel/Runtime/NKCanvas/src/NKCanvas/UI/NkGuiCanvasBackend.h:34, 41, 94, 120}
bool Init(nkentseu::renderer::NkIRenderer2D *renderer);
bool UploadFontGray8(uint32 texId, const uint8 *gray, int32 w, int32 h);
bool UploadImageRGBA(uint32 texId, const uint8 *rgba, int32 w, int32 h);
void Submit(const nkentseu::nkgui::NkGuiDrawList &dl, uint32 fbW, uint32 fbH);
\end{nkcode}

C'est tout. \texttt{Init} lui donne le renderer 2D ; les deux \texttt{Upload}
enregistrent une texture sous un identifiant (un atlas de police, une image) ;
\texttt{Submit} traduit une liste de commandes en appels
\texttt{NkIRenderer2D}. Le chapitre~4 montrera la séquence complète.

\subsection{\texttt{NkGuiRHIBackend} — le pont vers NKRHI}

L'alternative vit dans \nkfile{Integrations/NKGui/NkGuiRHIBackend.h}, compilée
dans une bibliothèque statique \texttt{NKGuiIntegration}. Son API est le miroir
de la précédente, avec un tampon de commandes en plus :

\begin{nkcode}{Integrations/NKGui/NkGuiRHIBackend.h:45-59}
bool Init(NkIDevice* device, NkRenderPassHandle renderPass, NkGraphicsApi api);
void Destroy();
void Submit(NkICommandBuffer* cmd, const NkGuiDrawList& dl, uint32 fbW, uint32 fbH);
bool UploadTextureRGBA8(uint32 texId, const uint8* data, int32 width, int32 height);
bool UploadTextureGray8(uint32 texId, const uint8* data, int32 width, int32 height);
bool RegisterTexture(uint32 texId, NkTextureHandle texture);
bool HasTexture(uint32 texId) const noexcept;
\end{nkcode}

La doctrine de répartition est écrite dans le fichier de build de l'intégration :

\begin{nkcode}{Integrations/NKGui/NKGuiIntegration.jenga:5-8}
Rend une UI NKGui (nkgui::NkGuiDrawList) via NKRHI/NKRenderer, sans NKCanvas.
Sert aux applications 2D/3D (app d'animation, moteur de jeu). NKCanvas reste
le backend de l'IDE (NKCode).
\end{nkcode}

Un point de conception intéressant de ce second pont : il enregistre aussi des
textures \emph{externes} (\texttt{RegisterTexture}). C'est ainsi qu'une
application 3D affiche le rendu de sa scène dans un panneau de son interface :
elle rend la scène dans une texture hors écran, l'enregistre sous un
\texttt{texId}, et NKGui la dessine comme une image ordinaire.

\subsection{Pourquoi cette séparation ?}

On pourrait trouver la construction alambiquée. Elle rapporte trois choses
concrètes :

\begin{enumerate}
  \item \textbf{On peut changer de moteur de rendu sans toucher un seul
        widget.} L'IDE utilise NkCanvas, l'application d'animation utilise NKRHI,
        et le code des boutons est le même. Mieux : NKEditorKit rend cette
        substitution explicite, avec une interface
        \texttt{NkIEditorRenderer} injectable
        (\nkfile{Engine/NKEditorKit/src/NKEditorKit/NkIEditorRenderer.h:7-15}).
  \item \textbf{NKGui est testable sans GPU.} Une liste de commandes est une
        structure de données : on peut la produire, l'inspecter et la comparer
        sans jamais ouvrir de fenêtre.
  \item \textbf{Les dépendances restent en arbre, pas en graphe.} NKGui ne
        connaît pas NkCanvas ; NkCanvas ne connaît pas NKGui ; l'application
        connaît les deux. Aucun cycle.
\end{enumerate}

Il y a un prix, et il faut le connaître : \textbf{c'est à l'application de faire
le raccordement}. Personne ne le fera à votre place. Si vous oubliez d'appeler
\texttt{Submit}, rien ne s'affiche et aucun message d'erreur n'apparaît. Si vous
oubliez d'uploader l'atlas de police, tous les textes sont invisibles et aucun
message d'erreur n'apparaît non plus. Le chapitre~4 revient longuement sur ces
silences.

\section{Le chemin complet d'un clic, de bout en bout}

Récapitulons en suivant un seul clic de souris, de l'appui jusqu'au pixel qui
change de couleur.

\begin{enumerate}
  \item L'utilisateur appuie. Le système envoie un message ; \texttt{NKWindow} le
        reçoit, \texttt{NKEvent} le transforme en
        \texttt{NkMouseButtonPressEvent}.
  \item La boucle appelle \texttt{PollEvent()} ; le \emph{callback} enregistré
        écrit \texttt{ctx.input.mouseDown[0] = true}.
  \item L'application appelle \texttt{ctx.BeginFrame(dt)}. NKGui en déduit la
        \emph{transition} : \texttt{mouseClicked[0] = mouseDown[0] \&\&
        !mousePrev[0]}.
  \item L'application appelle \texttt{Button(ctx, "OK")}. Le widget calcule son
        rectangle, demande son identité, teste le survol, voit le clic, et
        retourne \texttt{true} — au \emph{relâchement}, comme nous le verrons au
        chapitre~3. Au passage, il écrit deux rectangles et six sommets dans la
        liste de commandes.
  \item L'application appelle \texttt{ctx.EndFrame()}. NKGui fusionne les listes
        des fenêtres par ordre de profondeur, ferme les popups, remet à zéro la
        molette et le texte consommés.
  \item L'application appelle \texttt{target->Clear()}, puis soumet les deux
        listes au pont : \texttt{ctx.dl} d'abord, \texttt{ctx.dlOverlay}
        ensuite, chacune par \texttt{backend.Submit(liste, w, h)}. Le pont
        convertit chaque sommet, résout chaque \texttt{texId} en
        \texttt{NkTexture*}, applique chaque rectangle de découpe et appelle
        \texttt{DrawVertices}.
  \item NkCanvas accumule tout dans ses tampons, regroupe par texture, et à la
        fermeture de l'image envoie un appel de dessin par groupe.
  \item L'application appelle \texttt{target->Display()} : l'image est présentée.
\end{enumerate}

\begin{nkretenir}[L'invariant d'ordre]
\texttt{événements} $\rightarrow$ \texttt{ctx.BeginFrame(dt)} $\rightarrow$
\texttt{widgets} $\rightarrow$ \texttt{ctx.EndFrame()} $\rightarrow$
\texttt{Clear} $\rightarrow$ \texttt{Submit} $\times 2$ $\rightarrow$
\texttt{Display}.

Cet ordre n'est pas négociable et le chapitre~3 expliquera pourquoi
\texttt{BeginFrame} doit venir \emph{après} le pompage des événements.
\end{nkretenir}

\section{Deux mots sur les macros d'export}

En lisant les en-têtes, vous rencontrerez partout des macros du genre
\texttt{NKENTSEU\_NKGUI\_API} :

\begin{nkcode}{Kernel/Runtime/NKGui/src/NKGui/Widgets/NkGuiWidgets.h (forme générale)}
NKENTSEU_NKGUI_API void Text(NkGuiContext &ctx, const char *s) noexcept;
\end{nkcode}

Par défaut, le module est construit en \textbf{statique} et cette macro est
\textbf{vide} (\nkfile{Kernel/Runtime/NKGui/src/NKGui/NkGuiApi.h:32-40}). Quand
vous lisez une signature, effacez-la mentalement : il faut lire
\texttt{void Text(NkGuiContext \&ctx, const char *s) noexcept}. Il existe deux
variantes : \texttt{NKENTSEU\_NKGUI\_CLASS\_EXPORT} (classe entière) et
\texttt{NKENTSEU\_NKGUI\_API\_INLINE} (fonction \texttt{inline} exportée). Dans
la suite de ce cours, nous les omettons pour la lisibilité.

\begin{nkbilan}
Cinq étages, et une règle : chacun ignore celui du dessus. \textbf{NKWindow}
donne une fenêtre, \textbf{NKEvent} des événements typés, \textbf{NKCanvas} un
GPU abstrait --- distinct de NKRHI, et les deux ne se mélangent pas dans une
même fenêtre. \textbf{NkCanvas} dessine, \textbf{NKGui} produit une liste de
commandes et \emph{ne connaît aucun GPU} : c'est un pont qui traduit sa
\texttt{NkGuiDrawList}. Cette ignorance est ce qui permet à la même interface de
sortir sur NkCanvas ou sur NKRHI sans changer une ligne.
\end{nkbilan}

\section{Exercices}

\begin{nkdemo}[1 — Vérifier la carte soi-même]
Ouvrez les trois fichiers de build suivants et relevez, pour chacun, la liste
des modules dont il dépend :
\nkfile{Kernel/Runtime/NKGui/NKGui.jenga},
\nkfile{Kernel/Runtime/NKCanvas/NKCanvas.jenga},
\nkfile{Engine/NKEditorKit/NKEditorKit.jenga}.
Dessinez le graphe. Vérifiez qu'il n'y a aucun cycle, et identifiez le seul
module qui connaît à la fois NKGui et NKCanvas.
\end{nkdemo}

\begin{nklogique}[2 — Trouver le pont]
Ouvrez
\nkfile{Kernel/Runtime/NKCanvas/src/NKCanvas/UI/NkGuiCanvasBackend.h} et
comptez : combien de méthodes publiques ? Combien de lignes au total ? Comparez
avec les 4\,973 lignes de \texttt{NkGuiWidgets.cpp}. Que vous dit ce rapport sur
la quantité de travail nécessaire pour porter NKGui vers un nouveau moteur de
rendu ?
\end{nklogique}

\begin{nkpratique}[3 — Changer de backend à chaud]
Reprenez la démo \texttt{NkCanvasDemo} du dossier Sandbox, qui accepte un
argument de ligne de commande :
\begin{verbatim}
NkCanvasDemo.exe --backend=opengl
NkCanvasDemo.exe --backend=dx11
NkCanvasDemo.exe --backend=sw
\end{verbatim}
Lancez-la avec chacun des backends et notez ce que vous observez. Recoupez avec
le tableau d'état de la section 1.4 et avec \nkfile{logs/app.log}. Ce petit
exercice vous évitera plus tard de chercher un bug dans votre code alors que le
problème est dans le choix du backend.
\end{nkpratique}

\begin{nkquestions}
\begin{enumerate}
\item Citez les cinq étages de la pile et ce que chacun apporte.
\item Pourquoi NKCanvas et NKRenderer sont-ils exclusifs dans une même fenêtre ?
\item Que produit NKGui, et qu'est-ce qu'il ne fait \emph{pas} ?
\item À quoi sert un « pont » comme \texttt{NkGuiCanvasBackend} ?
\item Décrivez le trajet d'un clic, du système jusqu'au widget.
\end{enumerate}
\end{nkquestions}
